\documentclass[11pt,a4paper]{article}

\usepackage[margin=2.4cm]{geometry}
\usepackage[T1]{fontenc}
\usepackage[utf8]{inputenc}
\usepackage{lmodern}
\usepackage{microtype}
\usepackage{amsmath,amssymb,amsthm}
\usepackage{booktabs}
\usepackage{array}
\usepackage{algorithm}
\usepackage{algpseudocode}
\usepackage{pgfplots}
\usepackage{xcolor}
\usepackage{enumitem}
\usepackage[hidelinks]{hyperref}
\pgfplotsset{compat=1.17}
\usepgfplotslibrary{groupplots}
\setlength{\emergencystretch}{3em}

\newtheorem{theorem}{Theorem}
\newtheorem{lemma}[theorem]{Lemma}
\newtheorem{corollary}[theorem]{Corollary}
\newtheorem{proposition}[theorem]{Proposition}
\newtheorem{conjecture}[theorem]{Conjecture}
\theoremstyle{definition}
\newtheorem{definition}[theorem]{Definition}
\newtheorem{remark}[theorem]{Remark}

\newcommand{\NVerify}{280}
\newcommand{\NVerifyOk}{280}
\newcommand{\NMedium}{68}
\newcommand{\NMediumOk}{68}
\newcommand{\NScaling}{93}
\newcommand{\WindowMedianPct}{6.3}
\newcommand{\WindowMaxPct}{27.4}
\newcommand{\DecisionsSavedMean}{4.5}
\newcommand{\CertRate}{93}
\newcommand{\CertExactSmall}{348}
\newcommand{\NSmallAll}{348}

\newcommand{\VFiveInst}{6{,}000}
\newcommand{\VFiveTests}{201{,}422}
\newcommand{\VFiveFails}{5}
\newcommand{\ConeInst}{14{,}200}
\newcommand{\ConeTests}{429{,}742}
\newcommand{\ConeFails}{0}
\newcommand{\VFourInst}{14{,}000}
\newcommand{\VFourTests}{490{,}671}
\newcommand{\VFourFails}{0}
\newcommand{\VFourPertInst}{39{,}916}
\newcommand{\VFourPertTests}{876{,}002}
\newcommand{\VFourPertFails}{0}
\newcommand{\ConePertInst}{39{,}922}
\newcommand{\ConePertTests}{876{,}768}
\newcommand{\ConePertFails}{0}
\newcommand{\ConeAllTests}{1{,}306{,}510}

\newcommand{\TightFourChecked}{23{,}006}

\makeatletter
\newcommand{\inputrows}[1]{\@@input #1 }
\makeatother

\newcommand{\ProjectAuthors}{Project Team}
\newcommand{\opt}{r^{*}}
\newcommand{\Far}{\mathrm{far}}
\newcommand{\arc}[2]{P(#1,#2)}
\newcommand{\ang}{\angle}

\title{Exact $k$-Dispersion for Points in Convex Position:\\
  Delaunay Certificates, a Capped Dynamic Program,\\
  and Oriented Tight Pairs}
\author{\ProjectAuthors}
\date{\today}

\begin{document}
\maketitle

\begin{abstract}
Given $n$ points in convex position in the plane and an integer $k$, the
$k$-dispersion problem asks for $k$ of the points whose minimum pairwise
distance $\opt$ is as large as possible. We study exact algorithms and the
combinatorial structure of optimal solutions.

\emph{Algorithms.} We show that a set in convex position has all pairwise
distances at least $r$ if and only if it has a locally Delaunay triangulation
whose edges all have length at least $r$. This certificate gives a short
correctness proof for a triangulation dynamic program. Because a dispersion
decision only asks whether $k$ points fit, each chord of the program needs to
store only $k-2$ apex angles, compared exactly by integer in-circle tests.
The resulting decision procedure runs in $O(n^3\log k)$ time and $O(n^2k)$
space, and solves $k$-dispersion in $O(n^3\log k\log n)$ time, improving the
$O(n^3\log^2n)$ bound of Tkachenko and Wang (SWAT 2026) by a factor
$\log n/\log k$. A farthest-first greedy set and the hull perimeter bound
$\opt$ to the window $[g,\min(2g,\mathrm{per}(P)/k)]$.

\emph{Structure.} Call $\{q,a(q)\}$ a \emph{tight pair} when $a(q)$ is the
first point counterclockwise from $q$ at distance at least $r$; a clockwise
tight pair $\{b(q),q\}$ is defined symmetrically. We prove that whenever
$k$ points pairwise at distance $\ge r$ exist, some such $k$-set contains a
tight pair of one orientation. The proof combines a swap argument based on the
quadrilateral inequality with the fact that in every convex polygon some
vertex's nearest neighbour is adjacent to it. With a single prescribed
orientation the statement holds for $k=3$ but fails for $k=4$. We give an
explicit 8-point instance whose unique optimal set has no counterclockwise
tight pair. It refutes a conjecture that survived \ConeAllTests{} exact
randomised tests, and with it every search that grows solutions only from
such pairs. The oriented theorem yields an $O(n^2\log n)$ algorithm for $k=4$,
improving on the previous $O(n^2\log^2n)$. We also show that the arc
characterisation behind the near-linear algorithm for $k=3$ fails for every
$k\ge4$, and that natural alternating greedy rules are incomplete.

All algorithms are implemented with exact integer arithmetic and are
cross-checked against exhaustive search and an independent solver on a dataset
built from convex hulls of random point clouds.
\end{abstract}

\tableofcontents

%======================================================================
\section{Introduction}
%======================================================================

\paragraph{Problem.} Let $P$ be a set of $n$ points in the plane in
\emph{convex position}, i.e.\ every point of $P$ is a vertex of the convex hull
$\mathrm{conv}(P)$. For $2\le k\le n$, the \emph{$k$-dispersion} problem (also
called max-min dispersion or the maximally separated set problem) asks for
$S\subseteq P$, $|S|=k$, maximising $\delta(S)=\min_{p\ne q\in S}|pq|$. We
write $\opt=\opt_k(P)$. For general planar point sets the problem is
NP-hard~\cite{WK88}; the best exact algorithms run in
$n^{O(\sqrt k)}$~\cite{Akagi18} or $O(n^{k-1}\log n)$ time~\cite{CD25}. For
points in convex position it is polynomial. Dispersion models facility
placement in which facilities must be far apart, and convex position arises
when candidate sites lie on the boundary of a convex region.

\paragraph{Known results for convex position.} Table~\ref{tab:prior}
summarises the exact algorithms known to us (October 2026).

\begin{table}[h]
\centering
\caption{Exact algorithms for $k$-dispersion of $n$ points in convex position.}
\label{tab:prior}
\small
\begin{tabular}{@{}p{0.33\textwidth} p{0.3\textwidth} p{0.31\textwidth}@{}}
\toprule
Reference & Time & Remarks \\
\midrule
Singireddy, Basappa, Mitchell~\cite{SBM26} & $O(2^kn^2\log^2n)$ & bounded search tree, FPT in $k$\\
Singireddy, Basappa, Mitchell~\cite{SBM26} & $O(n^4k^2)$ & first polynomial algorithm \\
Tkachenko, Wang~\cite{TW25} & $O(n^{7/2}\log n)$; $O(n^{37/11}\log n)$ exp. & cuttings; randomised \\
Tkachenko, Wang~\cite{TW26} & $O(n^3\log^2n)$ & via $O(n^3\log n)$ max.\ independent set \\
Kobayashi et al.~\cite{Kob21} & $O(n^2)$ & $k=3$ \\
Tkachenko, Wang~\cite{TW25} & $O(n\log^2n)$; $O(n\log n)$ exp. & $k=3$ \\
\midrule
\textbf{This paper}, Thm.~\ref{thm:main} & $O(n^3\log k\log n)$ & decision $O(n^3\log k)$ \\
\textbf{This paper}, Thm.~\ref{thm:k4} & $O(n^2\log n)$ & $k=4$; decision $O(n^2)$ \\
\bottomrule
\end{tabular}
\end{table}

\paragraph{Contributions.}
\begin{enumerate}[leftmargin=*]
\item \textbf{Delaunay certificate} (Theorem~\ref{thm:cert}): a set in convex
  position is $r$-far (all pairwise distances $\ge r$) iff it has a
  triangulation in which every diagonal is locally Delaunay and every edge has
  length $\ge r$.
\item \textbf{Capped Delaunay DP} (Section~\ref{sec:dp}): the DP value of a
  chord, as a function of the angle constraint imposed by its parent, is a step
  function determined by $k-2$ apex points (Lemma~\ref{lem:thresholds}). This
  gives an exact decision procedure in $O(n^2k+T_r\log k)\subseteq O(n^3\log
  k)$ time and $O(n^2k)$ space using only integer orientation and in-circle
  tests, and an $O(n^3\log k\log n)$ algorithm (Theorem~\ref{thm:main}).
\item \textbf{Search window} (Lemma~\ref{lem:window}):
  $g\le\opt\le\min\{2g,\mathrm{per}(P)/k\}$.
\item \textbf{Oriented tight pairs} (Section~\ref{sec:tight}): every feasible
  instance has a solution containing a tight pair of some orientation
  (Theorem~\ref{thm:oriented}). With one fixed orientation this holds for $k=3$
  (Theorem~\ref{thm:k3}) and fails for $k=4$ (Proposition~\ref{prop:c1false}).
  As a consequence, $k=4$ can be solved in $O(n^2\log n)$ time
  (Theorem~\ref{thm:k4}).
\item \textbf{Limits of greedy structure} (Section~\ref{sec:greedy}): the arc
  characterisation used for $k=3$~\cite{Kob21,TW25} fails for $k\ge4$,
  alternating greedy rules are incomplete, and searches seeded only by
  counterclockwise tight pairs are incorrect. Every claim comes with an explicit
  counterexample, verified exhaustively in exact arithmetic.
\item \textbf{Implementation and experiments} (Section~\ref{sec:experiments}),
  including a \emph{certified} variant that runs an $O(n^2k)$ greedy and
  confirms its value with one exact decision.
\end{enumerate}

\paragraph{Positioning.} The bound $O(n^3\log k\log n)$ improves
$O(n^3\log^2n)$~\cite{TW26} by a factor $\log n/\log k$ and beats
$O(2^kn^2\log^2n)$~\cite{SBM26} once $2^k\gtrsim n/\log n$. It is the best
known deterministic bound roughly for $\log_2n\lesssim k\le n^{o(1)}$. For
$k=4$, Theorem~\ref{thm:k4} improves the previous best $O(n^2\log^2n)$ by a
$\log n$ factor. Our DP subproblems are those of~\cite{TW25} in substance. What
is new is the certificate-based proof, the $k$-capped representation, the
exact implementation, and the structural results. We could not access the full
texts of~\cite{TW26,SBM26} and compare with their stated bounds only.
Section~\ref{sec:greedy} shows that a seed-and-grow search whose seeds are only
counterclockwise tight pairs is incorrect for $k=4$. We do not know whether any
published algorithm relies on that property, and we offer the instance of
Proposition~\ref{prop:c1false} as a test case.

%======================================================================
\section{Preliminaries}
%======================================================================

\begin{table}[h]
\centering
\caption{Notation.}
\small
\begin{tabular}{@{}l p{0.72\textwidth}@{}}
\toprule
$P=(p_0,\dots,p_{n-1})$ & input points in strictly convex position, counterclockwise (ccw); indices mod $n$\\
$|pq|$ & Euclidean distance; all comparisons use squared distances (integers)\\
$\arc ij$ & the points strictly between $p_i$ and $p_j$ going ccw from $p_i$\\
$r$-far set & a set whose pairwise distances are all $\ge r$\\
$\delta(S)$, $\opt$ & minimum pairwise distance of $S$; $\opt=\max_{|S|=k}\delta(S)$\\
$a(q)$, $b(q)$ & first point ccw (resp.\ clockwise) from $q$ at distance $\ge r$\\
$\mathrm{succ}_S$, $\mathrm{pred}_S$ & next/previous point of $S$ in ccw order\\
$\ang abc$ & angle at $b$ in triangle $abc$, in $(0,\pi)$\\
$T_r$ & number of ccw triples $(i,l,j)$ that are pairwise $\ge r$\\
$\mathrm{per}(P)$ & perimeter of $\mathrm{conv}(P)$\\
\bottomrule
\end{tabular}
\end{table}

A \emph{triangulation} of a set $S$ in convex position is a maximal set of
non-crossing segments between points of $S$. Its edges are the sides of
$\mathrm{conv}(S)$ and the \emph{diagonals}. A diagonal $ab$ with incident
triangles $abc$ and $bad$ is \emph{locally Delaunay} if $d$ is not strictly
inside the circumcircle of $abc$, equivalently
$\ang acb+\ang adb\le\pi$. A triangulation is \emph{Delaunay} if no point of
$S$ is strictly inside the circumcircle of any triangle.

\paragraph{Exact predicates.} For integer coordinates of absolute value at
most $2^{29}$, the orientation and in-circle determinants fit in signed 128-bit
integers. All angle comparisons below reduce to them. For $x,y$ on the same
side of the line $ij$: $\ang ixj>\ang iyj$ iff $x$ is strictly inside the
circle through $i,y,j$. For a convex quadrilateral $i,x,l,j$:
$\ang ixl+\ang ijl>\pi$ iff $x$ is strictly inside the circle through
$i,l,j$.

\begin{lemma}[quadrilateral inequality]\label{lem:quad}
If $a,b,c,d$ are in strictly convex position in this cyclic order, then
$|ac|+|bd|>|ab|+|cd|$ and $|ac|+|bd|>|ad|+|bc|$.
\end{lemma}
\begin{proof}
The diagonals $ac$ and $bd$ cross at a point $o$ interior to both. By the
strict triangle inequality ($o$ lies on neither segment $ab$ nor $cd$),
$|ao|+|ob|>|ab|$ and $|co|+|od|>|cd|$. Adding gives
$|ac|+|bd|>|ab|+|cd|$. The second inequality is symmetric.
\end{proof}

%======================================================================
\section{A Delaunay certificate for \texorpdfstring{$r$}{r}-far sets}
%======================================================================

\begin{lemma}[Delaunay lemma~\cite{Del34,Ede01,BCKO08}]\label{lem:delaunay}
If every diagonal of a triangulation $T$ of a finite point set $S$ is locally
Delaunay, then $T$ is a Delaunay triangulation of $S$.
\end{lemma}

\begin{lemma}\label{lem:closest}
Let $T$ be a Delaunay triangulation of a set $S$ in strictly convex position,
$|S|\ge3$. Then $\delta(S)$ is attained by an edge of $T$.
\end{lemma}
\begin{proof}
Let $p,q$ be a closest pair and suppose $pq\notin T$. The triangles of $T$ at
$p$ partition the interior angle of $\mathrm{conv}(S)$ at $p$ into wedges. The
direction of $q$ from $p$ lies strictly inside the wedge of some triangle
$puv$ (not along an edge, by strict convexity), and $pq$ crosses $uv$. Let $C$
be the circumcircle of $puv$ and $A$ its arc from $u$ to $v$ not containing
$p$. Chords from $p$ inside the wedge end on $A$. Since $q$ is not strictly
inside $C$, the segment $pq$ leaves the closed disk at some $x\in A$, so
$|pq|\ge|px|$. Along $A$ the distance $|px|=2R\sin(\theta/2)$ (with $\theta$
the central angle) first increases and then decreases, so
$|px|\ge\min\{|pu|,|pv|\}$. Hence an edge of $T$ is at most as long as the
closest pair.
\end{proof}

\begin{theorem}[Delaunay certificate]\label{thm:cert}
Let $S$ be in strictly convex position, $|S|\ge3$, $r>0$. Then $S$ is $r$-far
iff $S$ has a triangulation $T$ with (a) every edge of length $\ge r$ and (b)
every diagonal locally Delaunay.
\end{theorem}
\begin{proof}
($\Leftarrow$) By Lemma~\ref{lem:delaunay} $T$ is Delaunay, and by
Lemma~\ref{lem:closest} $\delta(S)$ is the length of some edge of $T$, which
is $\ge r$. ($\Rightarrow$) A Delaunay triangulation of $S$ exists and
satisfies (b); (a) holds because $S$ is $r$-far.
\end{proof}

Conditions (a) and (b) are local, which is what makes a dynamic program
possible. The algorithm of~\cite{TW25} rests on the same structure, but its
correctness proof uses a seven-term distance inequality (their Lemma~14).

%======================================================================
\section{The capped Delaunay dynamic program}\label{sec:dp}
%======================================================================

Fix $r$. A chord $(i,j)$ is \emph{admissible} if $|p_ip_j|\ge r$; a point
$l\in\arc ij$ is an \emph{admissible apex} if $|p_ip_l|,|p_lp_j|\ge r$.

\begin{definition}\label{def:F}
For an admissible chord $(i,j)$, an \emph{$(i,j)$-configuration} is a pair
$(X,T)$ with $X\subseteq\arc ij$ and $T$ a triangulation of $X\cup\{p_i,p_j\}$
satisfying (a) and (b) of Theorem~\ref{thm:cert}. If $X\ne\emptyset$, its
\emph{apex} is the third vertex $l$ of the triangle of $T$ on $p_ip_j$. For
$\varphi\in(0,\pi]$ let $F(i,j\mid\varphi)$ be the maximum $|X|$ over
configurations with $X=\emptyset$ or $\ang p_ilp_j\le\varphi$.
\end{definition}

If $p_ip_j$ is later a diagonal whose other triangle is $p_jp_ip_t$, its local
Delaunay condition is exactly $\ang p_ilp_j\le\pi-\ang p_ip_tp_j$; the bound
$\varphi$ encodes this.

\begin{lemma}[recurrence]\label{lem:rec}
For every admissible chord $(i,j)$ and $\varphi\in(0,\pi]$:
\begin{multline*}
F(i,j\mid\varphi)=\max\Bigl(\{0\}\cup\bigl\{\,1+F(i,l\mid\pi-\ang p_ip_jp_l)+F(l,j\mid\pi-\ang p_lp_ip_j)\\
 \;:\; l \text{ admissible apex},\ \ang p_ilp_j\le\varphi\bigr\}\Bigr).
\end{multline*}
\end{lemma}
\begin{proof}
($\ge$) Combine optimal configurations for $(i,l)$ and $(l,j)$ with the
triangle $p_ip_lp_j$. The new edges have length $\ge r$. The new diagonals
$p_ip_l$ and $p_lp_j$ (when the children are nonempty) are locally Delaunay
exactly by the children's bounds. The apex is $l$ and satisfies the bound
$\varphi$. ($\le$) Removing the apex triangle of an optimal configuration
splits it into configurations for $(i,l)$ and $(l,j)$; condition (b) gives
their bounds.
\end{proof}

\begin{theorem}\label{thm:decision}
For $k\ge3$: $P$ has an $r$-far $k$-set iff some admissible chord has
$F(i,j\mid\pi)\ge k-2$.
\end{theorem}
\begin{proof}
($\Rightarrow$) If $S=\{q_1,\dots,q_k\}$ (ccw) is $r$-far, a Delaunay
triangulation of $S$ yields a $(q_1,q_k)$-configuration with
$|X|=k-2$ (Theorem~\ref{thm:cert}). ($\Leftarrow$) A configuration with
$|X|\ge k-2$ certifies, by Theorem~\ref{thm:cert}, an $r$-far set of size
$\ge k$.
\end{proof}

\paragraph{Thresholds.} For an admissible apex $l$ of $(i,j)$ let
$v(l)=1+F(i,l\mid\pi-\ang p_ip_jp_l)+F(l,j\mid\pi-\ang p_lp_ip_j)$, and for
$c\ge1$ let $\tau_c(i,j)=\min\{\ang p_ilp_j: v(l)\ge c\}$ (or $+\infty$).

\begin{lemma}\label{lem:thresholds}
$\tau_1\le\tau_2\le\cdots$ and $F(i,j\mid\varphi)=\max\{c\ge0:\tau_c\le\varphi\}$
(with $\tau_0=0$). With $K=k-2$ and $F_K=\min(F,K)$, the values
$\tau_1,\dots,\tau_K$ of every chord are determined by the values $F_K$ of
shorter chords.
\end{lemma}
\begin{proof}
The sets $\{l:v(l)\ge c\}$ shrink as $c$ grows, which gives monotonicity. By
Lemma~\ref{lem:rec}, $F(i,j\mid\varphi)\ge c$ iff some $l$ with
$\ang p_ilp_j\le\varphi$ has $v(l)\ge c$, i.e.\ iff $\tau_c\le\varphi$. Finally
$\min(K,v(l))=\min(K,1+F_K(\cdot)+F_K(\cdot))$, and $\tau_{c\le K}$ depends on
$v$ only through $\min(K,v)$.
\end{proof}

We store $\tau_c(i,j)$ as the apex $L_c(i,j)$ attaining it. The query
$F_K(i,l\mid\pi-\ang p_ip_jp_l)\ge c$ becomes ``$L_c(i,l)$ is not strictly
inside the circle through $p_i,p_l,p_j$'', which is monotone in $c$ and is
answered by binary search with $O(\log K)$ in-circle tests.

\begin{algorithm}[t]
\caption{\textsc{CappedDecision}$(P,k,r)$}\label{alg:capped}
\begin{algorithmic}[1]
\For{$g=2,\dots,n-1$; each admissible chord $(i,j)$ with $j=i+g$}
  \State $best[1..K]\gets\bot$
  \For{each admissible apex $l\in\arc ij$}
    \State $q_1\gets\#\{c: L_c(i,l)$ not strictly inside $\bigcirc(p_i,p_l,p_j)\}$;\;
           $q_2\gets\#\{c: L_c(l,j)$ not strictly inside $\bigcirc(p_i,p_l,p_j)\}$
    \State $v\gets\min(K,1+q_1+q_2)$; keep in $best[v]$ the apex of smaller angle
  \EndFor
  \State $L_c(i,j)\gets$ smallest-angle apex among $best[c..K]$, for $c=K,\dots,1$
\EndFor
\State \Return \textbf{true} iff some admissible chord has $L_K\ne\bot$
\end{algorithmic}
\end{algorithm}

\begin{theorem}\label{thm:capped}
Algorithm~\ref{alg:capped} decides whether $P$ has an $r$-far $k$-set in
$O(n^2k+T_r\log k)\subseteq O(n^3\log k)$ time and $O(n^2k)$ space, using
only exact integer predicates. A witness is recovered in $O(k\log k)$ time.
\end{theorem}
\begin{proof}
Correctness follows from Theorem~\ref{thm:decision} and
Lemma~\ref{lem:thresholds}. The inner loop runs $T_r$ times at $O(\log K)$
cost, and each chord costs $O(K)$ more. A witness is recovered by descending
from a chord with $L_K\ne\bot$: with demand $c$ at $(i,j)$, take
$l=L_c(i,j)$ and send $c_1=\min(q_1,c-1)$ to $(i,l)$ and $c-1-c_1$ to $(l,j)$.
\end{proof}

\begin{remark}
Without capping, a chord's step function can have $\Theta(n)$ steps, giving
$O(n^3\log n)$, the bound of~\cite{TW26}. Capping reduces the logarithm to
$\log k$ and the memory to $O(n^2k)$.
\end{remark}

%======================================================================
\section{Optimisation and the search window}
%======================================================================

\begin{lemma}[search window]\label{lem:window}
Let $G$ be the farthest-first greedy $k$-set and $g=\delta(G)$. Then
$g\le\opt\le\min\{2g,\mathrm{per}(P)/k\}$.
\end{lemma}
\begin{proof}
$g\le\opt$ since $G$ is feasible. The insertion distances of the greedy are
non-increasing, so every point of $P$ is within $g$ of the first $k-1$ greedy
points. Two points of an optimal set then share a nearest greedy point, so
$\opt\le2g$. For $k\ge3$, the convex hull of an optimal set is a convex
$k$-gon with all sides $\ge\opt$, contained in $\mathrm{conv}(P)$. Perimeter is
monotone under inclusion of convex sets, so $k\opt\le\mathrm{per}(P)$. For
$k=2$, $\opt$ is the diameter, which is at most $\mathrm{per}(P)/2$.
\end{proof}

\begin{theorem}\label{thm:main}
$k$-dispersion of $n$ points in convex position with integer coordinates in
$[-2^{29},2^{29}]$ is solvable exactly in $O(n^3\log k\log n)$ time and
$O(n^2k)$ space.
\end{theorem}
\begin{proof}
Binary search over the sorted distinct squared distances inside the window,
using Algorithm~\ref{alg:capped} as the decision procedure.
\end{proof}

%======================================================================
\section{Oriented tight pairs}\label{sec:tight}
%======================================================================

Fix $r$. For $q\in P$ with some point at distance $\ge r$, recall that $a(q)$
and $b(q)$ are the first such points counterclockwise and clockwise from $q$.
A pair $\{q,a(q)\}$ is a \emph{ccw-tight pair} and $\{b(q),q\}$ a
\emph{cw-tight pair}. In an $r$-far set $S$, a ccw-tight pair is consecutive:
$\mathrm{succ}_S(q)=a(q)$, since every point strictly between $q$ and $a(q)$ is
closer than $r$ to $q$.

\subsection{A swap lemma}

\begin{lemma}[swap]\label{lem:swap}
Let $S$ be $r$-far and $t\in S$.
\begin{enumerate}[label=(\roman*)]
\item If $s=\mathrm{pred}_S(t)$ is a nearest neighbour of $t$ in $S$, then
  $(S\setminus\{t\})\cup\{a(s)\}$ is an $r$-far $k$-set containing the
  ccw-tight pair $\{s,a(s)\}$.
\item If $s'=\mathrm{succ}_S(t)$ is a nearest neighbour of $t$ in $S$, then
  $(S\setminus\{t\})\cup\{b(s')\}$ is an $r$-far $k$-set containing the
  cw-tight pair $\{b(s'),s'\}$.
\end{enumerate}
\end{lemma}
\begin{proof}
(i) Since $|st|\ge r$, $w:=a(s)$ lies in $\arc st\cup\{t\}$. If $w=t$ there
is nothing to do. Otherwise $w$ is strictly between $s$ and $t$, and for each
$u\in S\setminus\{s,t\}$ the points $s,w,t,u$ are in strictly convex position
in this cyclic order. Lemma~\ref{lem:quad} gives
\[
  |st|+|wu|>|sw|+|tu|\ \ge\ r+|st|,
\]
using $|sw|\ge r$ (definition of $a$) and $|tu|\ge|ts|$ ($s$ is a nearest
neighbour of $t$). Hence $|wu|>r$. Also $|ws|\ge r$, so the new set is
$r$-far. (ii) is the mirror image: $w:=b(s')$ lies strictly between $t$ and
$s'$ (or equals $t$). For $u\in S\setminus\{t,s'\}$, the points $t,w,s',u$ are
in convex position, and $|ts'|+|wu|>|ws'|+|ut|\ge r+|ts'|$.
\end{proof}

\subsection{Nearest neighbours in convex polygons}

\begin{lemma}\label{lem:nn}
Let $Q$ be a convex polygon with $k\ge3$ vertices and choose for every vertex
$v$ a nearest neighbour $\nu(v)$. Then
\begin{enumerate}[label=(\roman*)]
\item the segments $v\nu(v)$ are pairwise non-crossing, and
\item for some vertex $v$, $\nu(v)$ is one of the two polygon neighbours of $v$.
\end{enumerate}
\end{lemma}
\begin{proof}
(i) Suppose $v\nu(v)$ and $y\nu(y)$ cross properly. Writing $v'=\nu(v)$ and
$y'=\nu(y)$, the four points are in convex position in the cyclic order
$v,y,v',y'$ (up to reversal), and $vy'$, $yv'$ are opposite sides.
Lemma~\ref{lem:quad} gives $|vv'|+|yy'|>|vy'|+|yv'|$. But $|vv'|\le|vy'|$ and
$|yy'|\le|yv'|$ by the choice of nearest neighbours, a contradiction.
(ii) Suppose every $v\nu(v)$ is a diagonal. Each diagonal $d$ leaves at least
one vertex strictly on each side; let $m(d)\ge1$ be the smaller of the two
counts. Choose $d=xy$ minimising $m(d)$, let $C$ be the side with $m(d)$
vertices, and pick $v\in C$. If $\nu(v)$ were on the other side of $xy$, the
segment $v\nu(v)$ would cross $xy$ properly, contradicting (i). So both
endpoints of the diagonal $v\nu(v)$ lie in $C\cup\{x,y\}$. One of its sides
then contains only vertices of $C\setminus\{v\}$, so
$m(v\nu(v))\le m(d)-1$, contradicting the choice of $d$.
\end{proof}

\subsection{The oriented tight-pair theorem}

\begin{theorem}\label{thm:oriented}
If $P$ has an $r$-far $k$-set ($k\ge2$), then it has one that contains a
ccw-tight pair or a cw-tight pair.
\end{theorem}
\begin{proof}
For $k=2$ take any $q$ with a point at distance $\ge r$; then $\{q,a(q)\}$
works. For $k\ge3$ let $S$ be $r$-far. By Lemma~\ref{lem:nn}(ii), applied to
the convex polygon $S$, some $t\in S$ has a polygon neighbour as its nearest
neighbour. Lemma~\ref{lem:swap}(i) or~(ii) then produces the required set.
\end{proof}

\begin{corollary}\label{cor:top}
In the DP of Section~\ref{sec:dp}, the top level may be restricted to the
$2n$ chords $(a(q),q)$ and $(q,b(q))$.
\end{corollary}

\subsection{One orientation: true for \texorpdfstring{$k=3$}{k=3}, false for \texorpdfstring{$k=4$}{k=4}}

Earlier versions of this work conjectured that one orientation always
suffices:

\smallskip\noindent\textbf{Property $C_1$.} \emph{If $P$ has an $r$-far
$k$-set, then it has one containing a ccw-tight pair.}

\begin{theorem}\label{thm:k3}
$C_1$ holds for $k\le3$ (and so does its mirror image).
\end{theorem}
\begin{proof}
$k=2$ is trivial. Let $S=\{x,y,z\}$ be $r$-far in ccw order. By
Lemma~\ref{lem:swap}(i) it suffices that for one of $(s,t,u)\in
\{(x,y,z),(y,z,x),(z,x,y)\}$ we have $|tu|\ge|ts|$. Otherwise
$|yz|<|xy|$, $|zx|<|yz|$ and $|xy|<|zx|$, i.e.\ $|xy|<|zx|<|yz|<|xy|$, which is
impossible.
\end{proof}

\begin{proposition}\label{prop:c1false}
$C_1$ fails for $k=4$, even at $r=\opt$. In the 8-point instance of
Table~\ref{tab:c1} (strictly convex, ccw order), the unique optimal 4-set is
$S=\{p_0,p_2,p_4,p_6\}$ with $(\opt)^2=555009055037$, attained by the diagonal
$p_2p_6$. At $r=\opt$ we have $a(p_0)=p_1$, $a(p_2)=p_3$, $a(p_4)=p_5$ and
$a(p_6)=p_7$, none of which is in $S$. So no $r$-far 4-set contains a ccw-tight
pair. (Every pair $\{b(q),q\}$ with $q\in S$ lies in $S$, as
Theorem~\ref{thm:oriented} predicts.)
\end{proposition}

\begin{table}[h]
\centering
\caption{The 8-point counterexample to $C_1$ for $k=4$ (integer
coordinates). Points with odd index are the ``blocking'' points $a(p_{2i})$.}
\label{tab:c1}
\small
\begin{tabular}{@{}c r r c c r r@{}}
\toprule
$i$ & $x$ & $y$ & \hspace{1em} & $i$ & $x$ & $y$ \\
\midrule
0 & $-778191$ & $-140663$ && 4 & $813287$ & $311458$ \\
1 & $-36313$ & $-275041$ && 5 & $72900$ & $437065$ \\
2 & $54127$ & $-289061$ && 6 & $-28204$ & $451365$ \\
3 & $643112$ & $175939$ && 7 & $-620451$ & $-15151$ \\
\bottomrule
\end{tabular}
\end{table}

The claims of Proposition~\ref{prop:c1false} were verified by exhaustive
enumeration in exact integer arithmetic by two independent programs.

\paragraph{How the counterexample was found.} Randomised exact tests of $C_1$
on \ConeInst{}+\ConePertInst{} instances (\ConeAllTests{} decisions, including
targeted perturbations of every other counterexample in this paper) found no
failure. The construction instead came from analysing the swap. The swap at
$s=\mathrm{pred}_S(t)$ fails only if $a(s)$, which lies in the \emph{cap}
beyond the edge $st$ of $\mathrm{conv}(S)$, is closer than $r$ to some other
vertex of $S$. By Lemma~\ref{lem:quad} this needs $|tu|<|ts|$ for that vertex.
We therefore searched, by hill climbing over quadrilaterals, for a convex
$S$ in which every cap contains such a point. We then placed one point in each
cap, rounded to integers, and checked the result exhaustively. The relevant
configurations have small measure, which explains why random sampling missed
them; this is a cautionary example for experimental evidence on such
questions.

\subsection{An \texorpdfstring{$O(n^2\log n)$}{O(n² log n)} algorithm for \texorpdfstring{$k=4$}{k=4}}

\begin{theorem}\label{thm:k4}
For $k=4$ the decision problem can be solved in $O(n^2)$ time, and
$4$-dispersion of $n$ points in convex position in $O(n^2\log n)$ time.
\end{theorem}
\begin{proof}
By Theorem~\ref{thm:oriented}, an $r$-far 4-set exists iff for some $q$ and
some orientation, with $(x,y)=(q,a(q))$ or $(b(q),q)$, there are two points
$u,u'$ in the arc of $P$ opposite to the pair, both at distance $\ge r$ from
$x$ and $y$ and with $|uu'|\ge r$. (Given a solution with a ccw-tight pair
$\{q,a(q)\}$, its other two points lie in $\arc{a(q)}{q}$; symmetrically for
cw.) For fixed $q$ and orientation, the set $F$ of points of that arc far from
both $x$ and $y$ is computed in $O(n)$ time. $F$ is in convex position in
ccw order, so its farthest pair is found in $O(|F|)$ time by rotating
calipers~\cite{Sha78}. Computing $a(\cdot)$ and $b(\cdot)$ takes $O(n^2)$
time, and $2n$ choices of $(q,\text{orientation})$ give $O(n^2)$ in total.
Sorting the candidate distances takes $O(n^2\log n)$, and the binary search
makes $O(\log n)$ decisions.
\end{proof}

\begin{corollary}\label{cor:k5}
For $k=5$, replacing the farthest-pair query by the $O(m\log m)$ decision
procedure for three points of~\cite{TW25} (Theorem~5 there) gives an
$O(n^2\log n)$ decision procedure and an $O(n^2\log^2n)$ algorithm.
\end{corollary}

%======================================================================
\section{Limits of greedy structure}\label{sec:greedy}
%======================================================================

\paragraph{The arc characterisation.} For $k=3$ there is an $O(n\log n)$
decision procedure~\cite{TW25}, based on the characterisation (their Lemma~16,
from~\cite{Kob21}): with $A(i)$ the ccw arc $P[a(p_i),b(p_i)]$,
\[
(G_k)\qquad P\text{ has an $r$-far $k$-set}\iff\exists i:\ A(i)\text{ contains an $r$-far $(k-1)$-set}.
\]
$G_3$ is true. $G_4$ is false: Table~\ref{tab:counter} lists a 6-point
instance with no $r$-far 4-set in which an arc $A(i)$ contains an $r$-far
triple. Our random tests also found failures of $G_5$ and $G_6$.

\paragraph{Seeded greedy search.} A seed-and-grow search keeps a ccw and a
clockwise frontier and repeatedly adds the first point beyond a frontier that
is $r$-far from all chosen points. We tested the following variants:
\begin{itemize}[leftmargin=*]
\item[$V_1$/$V_3$:] seed $\{q,a(q),b(q)\}$, alternating sides / all side patterns;
\item[$V_5$/$V_4$:] seed $\{q,a(q)\}$, alternating / all patterns;
\item[$V_4^{\pm}$:] seeds $\{q,a(q)\}$ and $\{b(q),q\}$, all patterns.
\end{itemize}

\begin{proposition}\label{prop:greedy}
$V_1$ and $V_3$ fail for $k=4$, and $V_5$ fails for $k=6$
(Table~\ref{tab:counter}). $V_4$ fails for $k=4$: it implies $C_1$, which
fails by Proposition~\ref{prop:c1false}.
\end{proposition}

\begin{table}[h]
\centering
\caption{Counterexamples (points in ccw order, $0$-based indices). For $G_4$
no $r$-far 4-set exists; in the other rows an $r$-far $k$-set exists (one is
shown) but the property fails.}
\label{tab:counter}
\small
\begin{tabular}{@{}l c c r l@{}}
\toprule
Property & $n$ & $k$ & $r^2$ & an $r$-far $k$-set \\
\midrule
\inputrows{generated/counterexamples.tex}
\bottomrule
\end{tabular}
\end{table}

All counterexamples are thin convex polygons. In such polygons, points that
are adjacent in space can be far apart in the cyclic order, so a search that
commits early on one side can block the other.

\begin{conjecture}\label{conj:v4pm}
$V_4^{\pm}$ is complete: whenever $P$ has an $r$-far $k$-set, the search
seeded with some oriented tight pair and some side pattern reaches $k$ points.
\end{conjecture}

The conjecture holds for $k\le3$ by Theorem~\ref{thm:oriented}. It held in all
\VFourPMTests{} exact tests (\VFourPMInst{} instances, including perturbations
of all counterexamples), and on every instance of Table~\ref{tab:counter}. If
true, it gives an $O(2^kn^2k\log n)$ algorithm: the search tree has at most
$2^{k-1}$ nodes per seed and each node scans $O(n)$ points against $\le k$
chosen ones. Given Proposition~\ref{prop:c1false}, we regard this evidence
with caution.

\paragraph{A certified algorithm.} Seeded greedies are sound: every set they
output is $r$-far. Algorithm \textsc{Certified} binary-searches the window with
the greedy (seeds of both orientations, alternating) as predicate, obtaining
$\rho\le\opt$. It then runs one exact decision at the successor of $\rho$. If
that fails, $\rho=\opt$; otherwise the exact binary search continues above
$\rho$. The output is always exact, and the time is
$O(n^2k\log n+n^3\log k)$ whenever the greedy finds $\opt$.

%======================================================================
\section{Experiments}\label{sec:experiments}
%======================================================================

\paragraph{Dataset (random points $\to$ convex hull).} Instances are generated
by sampling $N$ random integer points (coordinates up to $10^6$) from a
distribution, taking the strictly convex hull, and keeping $n$ of its
vertices. Any subset of a set in convex position is in convex position. The
distributions are uniform square and Gaussian (small hulls, so $n\le30$ and
$n\le14$), uniform disk, thin annulus, thin $1{:}20$ ellipse, points on a
circle, and \emph{cocircular} points: all 324 lattice points on
$x^2+y^2=32045^2$, the most degenerate input for in-circle tests. There are
three splits: \emph{verify} ($n\le14$, \NVerify{} instances), \emph{medium}
($n\le40$, \NMedium) and \emph{scaling} ($n\le200$, \NScaling).

\paragraph{Correctness.} Each instance was solved by the capped DP, the
certified algorithm, the Tkachenko--Wang DP (direct implementation),
exhaustive search and, on \emph{verify}, an independent FMM-based clique
solver. Every returned set was re-checked independently. The $k=4$ algorithm
of Theorem~\ref{thm:k4} agreed with the capped DP on \TightFourChecked{}
random instances and on all $k=4$ counterexamples.

\begin{table}[h]
\centering
\caption{Optimum agrees with exhaustive search and witness valid, per
distribution: \NVerifyOk/\NVerify{} (verify) and \NMediumOk/\NMedium{}
(medium) instances fully correct.}
\begin{tabular}{@{}l c c c c c@{}}
\toprule
\multicolumn{6}{@{}l}{\emph{verify} ($n\le14$)} \\
Distribution & inst. & capped & certified & TW & FMM \\
\midrule
\inputrows{generated/verify_by_dist.tex}
\midrule
\multicolumn{6}{@{}l}{\emph{medium} ($20\le n\le40$)} \\
Distribution & inst. & capped & certified & TW & \\
\midrule
\inputrows{generated/medium_by_dist.tex}
\bottomrule
\end{tabular}
\end{table}

\paragraph{Running time.} Table~\ref{tab:scaling} and
Figure~\ref{fig:scaling} report median in-process solve times on one core.
The window kept a median of \WindowMedianPct\% (max \WindowMaxPct\%) of the
candidate distances and saved \DecisionsSavedMean{} exact decisions on
average. The certified algorithm was certified by its greedy on
\CertRate/\NScaling{} scaling instances.

\begin{table}[h]
\centering
\caption{Median solve time [ms] and exact decisions on the scaling split
(TW only for $n\le100$).}
\label{tab:scaling}
\small
\begin{tabular}{@{}c c r r r r c c c@{}}
\toprule
$k$ & $n$ & TW & capped & certified & TW/capped & dec.\ (no window) & dec.\ (window) & dec.\ (cert.) \\
\midrule
\inputrows{generated/scaling_table.tex}
\bottomrule
\end{tabular}
\end{table}

\begin{figure}[h]
\centering
\begin{tikzpicture}
\begin{groupplot}[group style={group size=4 by 1, horizontal sep=0.9cm},
  width=0.27\textwidth, height=0.28\textwidth, ymode=log, xlabel={$n$},
  unbounded coords=jump, grid=major, title style={font=\small},
  tick label style={font=\scriptsize},
  legend style={font=\scriptsize, at={(0.03,0.97)}, anchor=north west}]
\nextgroupplot[title={$k=4$}, ylabel={time [ms]}]
 \addplot[mark=triangle*, black] table[x=n, y=tw] {generated/scaling_k4.dat};
 \addplot[mark=*, blue] table[x=n, y=capped] {generated/scaling_k4.dat};
 \addplot[mark=square*, red] table[x=n, y=certified] {generated/scaling_k4.dat};
 \legend{TW DP, capped, certified}
\nextgroupplot[title={$k=8$}]
 \addplot[mark=triangle*, black] table[x=n, y=tw] {generated/scaling_k8.dat};
 \addplot[mark=*, blue] table[x=n, y=capped] {generated/scaling_k8.dat};
 \addplot[mark=square*, red] table[x=n, y=certified] {generated/scaling_k8.dat};
\nextgroupplot[title={$k=16$}]
 \addplot[mark=triangle*, black] table[x=n, y=tw] {generated/scaling_k16.dat};
 \addplot[mark=*, blue] table[x=n, y=capped] {generated/scaling_k16.dat};
 \addplot[mark=square*, red] table[x=n, y=certified] {generated/scaling_k16.dat};
\nextgroupplot[title={$k=4$, large $n$}, xmode=log]
 \addplot[mark=*, blue] table[x=n, y=capped] {generated/k4_scaling.dat};
 \addplot[mark=diamond*, teal] table[x=n, y=tight4] {generated/k4_scaling.dat};
 \legend{capped, Thm.~\ref{thm:k4}}
\end{groupplot}
\end{tikzpicture}
\caption{Median solve time vs.\ $n$ (log scale). The right panel compares the
$k=4$ algorithm of Theorem~\ref{thm:k4} with the capped DP on random convex
polygons, two-chain shapes and circles.}
\label{fig:scaling}
\end{figure}

\paragraph{Structural tests.} Table~\ref{tab:structure} summarises the
exhaustive tests of structural properties.

\begin{table}[h]
\centering
\caption{Exact tests of structural properties: counts of $(P,k,r)$ tested
and failures found by random search. ``Status'' is the proven status.}
\label{tab:structure}
\small
\begin{tabular}{@{}l c r r l@{}}
\toprule
Property & $k$ & tests & failures & status \\
\midrule
$C_1$ (random hunts, $n\le30$) & 3--12 & \ConeTests & \ConeFails & false for $k=4$ (Prop.~\ref{prop:c1false}) \\
$C_1$ (perturbations) & 3--10 & \ConePertTests & \ConePertFails & \\
$V_4$ (random hunts, $n\le26$) & 3--12 & \VFourTests & \VFourFails & false for $k=4$ (Prop.~\ref{prop:greedy}) \\
$V_4$ (perturbations) & 3--10 & \VFourPertTests & \VFourPertFails & \\
$V_5$ (random hunts, $n\le22$) & 3--9 & \VFiveTests & \VFiveFails & false for $k=6$ \\
$V_4^{\pm}$ (hunts and perturbations) & 3--10 & \VFourPMTests & \VFourPMFails & open (Conj.~\ref{conj:v4pm}) \\
\bottomrule
\end{tabular}
\end{table}

%======================================================================
\section{Conclusion and open problems}
%======================================================================

We gave a simple exact $O(n^3\log k\log n)$ algorithm for $k$-dispersion in
convex position, an $O(n^2\log n)$ algorithm for $k=4$, and a structural
theory of tight pairs: some optimal set always contains a tight pair of
\emph{some} orientation, but not necessarily of a prescribed one once $k\ge4$.
Open questions:
\begin{enumerate}[leftmargin=*]
\item Does $C_1$ fail for every $k\ge4$? Our lifting experiments ($k=5$) found
  no counterexample.
\item Prove or refute Conjecture~\ref{conj:v4pm}. A proof would give an
  $O(2^kn^2k\log n)$ algorithm with a self-contained correctness argument.
\item Can $k=4$ be solved in $o(n^2)$ time, matching the near-linear bound for
  $k=3$? Is there an $o(n^3)$ algorithm for every fixed $k$?
\item Can the $\log n$ factor of the binary search be removed?
\end{enumerate}

\paragraph{Reproducibility.} Code, data and scripts are in \texttt{convex/}:
the solver (\texttt{convex\_dispersion.cpp}), dataset generation, the
cross-verification harness, structural experiments, counterexample search,
minimisation and independent re-verification
(\texttt{experiments/check\_counterexamples.py}).

\begin{thebibliography}{99}
\bibitem{Akagi18} T.~Akagi et al., Exact algorithms for the max-min dispersion
  problem, \emph{FAW 2018}, LNCS 10823, 263--272.
\bibitem{BCKO08} M.~de Berg, O.~Cheong, M.~van Kreveld, M.~Overmars,
  \emph{Computational Geometry: Algorithms and Applications}, 3rd ed.,
  Springer, 2008.
\bibitem{CD25} K.~Chen, A.~Dumitrescu, A couple of simple algorithms for
  $k$-dispersion, arXiv:2511.00692, 2025.
\bibitem{Del34} B.~Delaunay, Sur la sph\`ere vide, \emph{Bull.\ Acad.\ Sci.\
  URSS} 6 (1934), 793--800.
\bibitem{Ede01} H.~Edelsbrunner, \emph{Geometry and Topology for Mesh
  Generation}, Cambridge University Press, 2001.
\bibitem{Kob21} Y.~Kobayashi, S.-I.~Nakano, K.~Uchizawa, T.~Uno, Y.~Yamaguchi,
  K.~Yamanaka, An $O(n^2)$-time algorithm for computing a max-min 3-dispersion
  on a point set in convex position, \emph{IEICE Trans.\ Inf.\ \& Syst.}
  E105-D(3), 2022.
\bibitem{Sha78} M.~I.~Shamos, \emph{Computational Geometry}, PhD thesis, Yale
  University, 1978.
\bibitem{SBM26} V.~R.~Singireddy, M.~Basappa, J.~S.~B.~Mitchell, Algorithms for
  $k$-dispersion for points in convex position in the plane, \emph{Discrete
  Applied Mathematics}, 2026 (conference version CALDAM 2023).
\bibitem{TW25} A.~Tkachenko, H.~Wang, Dominating set, independent set,
  discrete $k$-center, dispersion, and related problems for planar points in
  convex position, \emph{STACS 2025}, LIPIcs 327, 73; arXiv:2501.00207.
\bibitem{TW26} A.~Tkachenko, H.~Wang, Maximum independent sets in disk graphs
  with disks in convex position, \emph{SWAT 2026}, LIPIcs 370, 40;
  arXiv:2604.10828.
\bibitem{WK88} D.-W.~Wang, Y.-S.~Kuo, A study on two geometric location
  problems, \emph{Inform.\ Process.\ Lett.} 28(6) (1988), 281--286.
\end{thebibliography}

\end{document}
